%% ============================================================
%%  Blog Post — Deep Learning Assignment 2
%%  Paste into Overleaf. Upload the three images:
%%    dataset_samples.png
%%    dataset_comparison.png
%%    rotation_accuracy.png
%% ============================================================
\documentclass[12pt,a4paper]{article}

% ── packages ─────────────────────────────────────────────────
\usepackage[top=2.5cm, bottom=2.5cm, left=3cm, right=3cm]{geometry}
\usepackage{fontenc}
\usepackage{inputenc}
\usepackage{lmodern}
\usepackage{microtype}
\usepackage{graphicx}
\usepackage{xcolor}
\usepackage{booktabs}
\usepackage{tabularx}
\usepackage{array}
\usepackage{amsmath}
\usepackage{enumitem}
\usepackage{hyperref}
\usepackage{parskip}
\usepackage{titlesec}
\usepackage{mdframed}
\usepackage{tcolorbox}
\usepackage{caption}
\usepackage{fancyhdr}
\usepackage{soul}

% ── colours ──────────────────────────────────────────────────
\definecolor{blogblue}{HTML}{1a73e8}
\definecolor{bloggray}{HTML}{5f6368}
\definecolor{blogyellow}{HTML}{fff3cd}
\definecolor{blogyellowborder}{HTML}{e0a800}
\definecolor{bloggreen}{HTML}{1e8449}
\definecolor{blogred}{HTML}{c0392b}
\definecolor{bloglightgray}{HTML}{f8f9fa}
\definecolor{blogborder}{HTML}{dee2e6}
\definecolor{codebg}{HTML}{f4f4f4}

% ── hyperlinks ────────────────────────────────────────────────
\hypersetup{
  colorlinks=true,
  linkcolor=blogblue,
  urlcolor=blogblue,
  citecolor=blogblue
}

% ── section styling (blog-like) ───────────────────────────────
\titleformat{\section}
  {\Large\bfseries\color{black}}
  {}
  {0em}
  {}
  [\vspace{-0.3em}\textcolor{blogblue}{\rule{\linewidth}{2pt}}\vspace{0.3em}]

\titleformat{\subsection}
  {\large\bfseries\color{black}}
  {}
  {0em}
  {}

\titlespacing{\section}{0pt}{1.5em}{0.5em}
\titlespacing{\subsection}{0pt}{1em}{0.3em}

% ── caption styling ───────────────────────────────────────────
\captionsetup{
  font=small,
  labelfont={bf,color=bloggray},
  textfont={color=bloggray},
  justification=centering
}

% ── header / footer ───────────────────────────────────────────
\pagestyle{fancy}
\fancyhf{}
\renewcommand{\headrulewidth}{0pt}
\fancyhead[L]{\small\color{bloggray}Deep Learning — Assignment 2}
\fancyhead[R]{\small\color{bloggray}Control Dataset Blog Post}
\fancyfoot[C]{\small\color{bloggray}\thepage}

% ── callout box ───────────────────────────────────────────────
\tcbuselibrary{breakable, skins}

\newtcolorbox{hypothesisbox}{
  colback=blogyellow,
  colframe=blogyellowborder,
  boxrule=1pt,
  arc=4pt,
  left=10pt, right=10pt, top=8pt, bottom=8pt,
  breakable
}

\newtcolorbox{findingbox}{
  colback=bloglightgray,
  colframe=blogborder,
  boxrule=1pt,
  arc=4pt,
  left=10pt, right=10pt, top=8pt, bottom=8pt,
  breakable
}

\newtcolorbox{codebox}{
  colback=codebg,
  colframe=blogborder,
  boxrule=0.5pt,
  arc=3pt,
  left=10pt, right=10pt, top=6pt, bottom=6pt,
  fontupper=\ttfamily\small,
  breakable
}

% ── story-point badge ─────────────────────────────────────────
\newcommand{\storypoint}[2]{%
  \noindent
  \begin{tcolorbox}[
    colback=blogblue!8,
    colframe=blogblue,
    boxrule=0.8pt,
    arc=4pt,
    left=8pt, right=8pt, top=5pt, bottom=5pt,
    title={\textbf{\textcolor{white}{#1}}},
    coltitle=white,
    attach boxed title to top left={yshift=-2mm, xshift=6pt},
    boxed title style={colback=blogblue, arc=2pt}
  ]
  #2
  \end{tcolorbox}
  \vspace{4pt}
}

% ─────────────────────────────────────────────────────────────
\begin{document}

% ── title block ───────────────────────────────────────────────
\begin{center}
  {\fontsize{22}{28}\selectfont\bfseries
    Does Multi-View Projection Make SimpleView\\[4pt]
    More Rotation-Robust?}\\[12pt]
  {\large\color{bloggray} A Controlled Experiment}\\[8pt]
  \textcolor{blogborder}{\rule{0.5\linewidth}{1pt}}\\[8pt]
  {\color{bloggray}
    \textbf{Deep Learning} \quad|\quad
    Assignment 2 --- Toy Problem / Control Dataset \quad|\quad
    \today}
\end{center}

\vspace{1.5em}

% ─────────────────────────────────────────────────────────────
\section{Motivation and Hypothesis}

Automatic 3D object classification is essential for autonomous driving and robotic manipulation, where LiDAR sensors produce \textbf{point clouds} --- unordered sets of 3D coordinates. A key practical challenge: sensors don't guarantee that objects appear in a canonical orientation. A robust classifier must recognise objects regardless of pose.

Two families of methods tackle this differently.

\textbf{Point-based methods} operate directly on raw 3D coordinates.
\href{https://arxiv.org/abs/1612.00593}{PointNet (Qi et al., 2017)} applies shared MLPs per point and pools a global descriptor. It includes a \emph{T-Net} --- a sub-network predicting a $3\times3$ transformation matrix applied to the input before classification, an explicit mechanism to handle rotation.
\href{https://arxiv.org/abs/1801.07829}{DGCNN (Wang et al., 2019)} constructs a local $k$-NN graph and learns edge features between neighbours.

\textbf{Projection-based methods} convert point clouds to 2D images first.
\href{https://arxiv.org/abs/2106.05304}{SimpleView (Goyal et al., 2021)} renders each point cloud as \textbf{six orthogonal depth images} ($\pm X, \pm Y, \pm Z$ at $128\times128$), processes each with a compact ResNet-18, and classifies the concatenation of the six feature vectors.

On ModelNet40, SimpleView matches or outperforms point-based methods with fewer parameters. On \textbf{ScanObjectNN} --- a real-world dataset where objects are \emph{not} axis-aligned and are surrounded by background clutter --- SimpleView achieves \textbf{79.5\% accuracy}, outperforming DGCNN (78.1\%).

\medskip
\noindent\textbf{But ScanObjectNN confounds multiple variables at once.} Objects are non-axis-aligned \emph{and} partially occluded \emph{and} surrounded by background points. It is impossible to tell from these results alone whether SimpleView's advantage comes from:

\begin{itemize}[noitemsep, leftmargin=2em]
  \item \textbf{Rotation robustness} --- better handling of pose variation?
  \item \textbf{Noise/occlusion robustness} --- better handling of clutter?
  \item Some other factor entirely?
\end{itemize}

\medskip
This matters because the two families handle rotation \emph{fundamentally differently}. SimpleView projects to six \emph{fixed world-space} directions --- if an object rotates, each depth image changes because a different face is visible from each fixed camera. PointNet's T-Net, by contrast, attempts to explicitly canonicalise the input before processing.

\begin{hypothesisbox}
\textbf{Hypothesis:} If SimpleView's advantage on ScanObjectNN comes from rotation robustness, it should degrade \emph{more gracefully} than PointNet and DGCNN when objects are rotated at test time --- seeing all six faces simultaneously should help.

\medskip
\textbf{This is directly testable} with a fully controlled experiment where rotation is the \emph{only} variable (van Gemert RP9: change one variable at a time).
\end{hypothesisbox}

% ─────────────────────────────────────────────────────────────
\section{Dataset Description and Examples}

To isolate rotation as the single controlled variable, we generate a synthetic dataset of \textbf{five geometrically distinct 3D shape classes}:

\begin{center}
\begin{tabular}{ll}
\toprule
\textbf{Class} & \textbf{Sampling method} \\
\midrule
Cube     & Uniform over all six faces \\
Sphere   & Normalised Gaussian vectors (exactly uniform surface) \\
Cylinder & Area-proportional: lateral surface $+$ two circular caps \\
Cone     & Area-proportional: lateral surface $+$ base disk \\
Torus    & Rejection sampling, weight $(R + r\cos\theta)/(R+r)$ \\
\bottomrule
\end{tabular}
\end{center}

Every point cloud contains exactly \textbf{1,024 points}, normalised to zero-mean and unit maximum radius.

\medskip
\noindent\textbf{What the dataset does NOT include:} no noise, no background clutter, no occlusion, no class imbalance, no difference in point density. The only thing that changes between training and any test condition is the \textbf{rotation angle applied to each point cloud}. This strict isolation is what makes it a valid control experiment.

\subsection*{Dataset Splits}

\begin{center}
\begin{tabular}{lll}
\toprule
\textbf{Split} & \textbf{Rotation applied} & \textbf{Samples} \\
\midrule
\texttt{train.h5} & None (canonical orientation) & 2,000 \\
\texttt{test\_rot000.h5} & None (aligned) & 500 \\
\texttt{test\_rot030.h5} \ldots \texttt{test\_rot180.h5} & Y-axis $\theta^\circ$ in steps of $30^\circ$ & 500 each \\
\texttt{test\_so3\_t0.h5} \ldots \texttt{test\_so3\_t2.h5} & Random SO(3) rotation & 500 each (3 seeds) \\
\bottomrule
\end{tabular}
\end{center}

\subsection*{Visualisation}

Figure~\ref{fig:samples} shows all five classes at $0^\circ$ (aligned) and $90^\circ$ Y-axis rotation. Colour encodes height ($z$-value).

\begin{figure}[h]
  \centering
  \includegraphics[width=\linewidth]{dataset_samples.png}
  \caption{Five synthetic point-cloud classes at $0^\circ$ (top) and $90^\circ$ Y-axis rotation (bottom).}
  \label{fig:samples}
\end{figure}

Figure~\ref{fig:comparison} reveals the key structural property of the dataset that drives all results: three classes are rotationally symmetric (green), two are not (red).

\begin{figure}[h]
  \centering
  \includegraphics[width=\linewidth]{dataset_comparison.png}
  \caption{%
    \textcolor{bloggreen}{\textbf{Green (symmetric):}} sphere, torus, cube --- look identical at $0^\circ$ and $90^\circ$, always classified correctly.
    \textcolor{blogred}{\textbf{Red (asymmetric):}} cylinder, cone --- change appearance substantially at $90^\circ$, causing all models to fail.
    Floor = $3/5 = \mathbf{60\%}$.
  }
  \label{fig:comparison}
\end{figure}

% ─────────────────────────────────────────────────────────────
\section{Method of Generation}

All generation code is scripted end-to-end for full reproducibility (van Gemert RP11). No manual steps.

\subsection*{Shape Samplers}

Each class uses a custom, area-correct sampler to ensure uniform surface coverage:

\begin{itemize}[leftmargin=2em]

\item \textbf{Cube} --- Points drawn uniformly on each of the 6 faces with per-face counts proportional to face area (all equal for a unit cube).

\item \textbf{Sphere} --- Sample from $\mathcal{N}(0, I_3)$, normalise to unit sphere. This gives exactly uniform surface coverage.

\item \textbf{Cylinder} (radius$=1$, height$=2$) --- Lateral area $= 4\pi$, each cap area $= \pi$. Lateral: uniform $\theta \in [0,2\pi]$, uniform $z \in [-1,1]$. Caps: $r = \sqrt{u},\ u \sim \mathrm{Uniform}[0,1]$ (area-correct disk sampling).

\item \textbf{Cone} (apex at $z=+1$, base radius$=1$) --- Slant $= \sqrt{5}$. Lateral: $r \sim \sqrt{u}$ for area-correct sampling, $z = 1-2r$, $\theta \sim \mathrm{Uniform}[0,2\pi]$.

\item \textbf{Torus} ($R=0.65$, $r=0.25$) --- Rejection sampling: draw $(\theta, \phi) \sim \mathrm{Uniform}[0,2\pi]^2$, accept with probability $(R+r\cos\theta)/(R+r)$ for uniform surface density.

\end{itemize}

\subsection*{Normalisation}

Every point cloud is zero-centred (subtract mean) and scaled to the unit sphere (divide by max radius). All classes occupy the same spatial scale.

\subsection*{Rotation}

\begin{itemize}[noitemsep, leftmargin=2em]
  \item \textbf{Y-axis:} $3\times3$ rotation matrix $R_y(\theta)$ applied to all points.
  \item \textbf{SO(3):} \texttt{scipy.spatial.transform.Rotation.random()} for uniformly distributed random orientations. Three independent seeds for stability.
\end{itemize}

\subsection*{The Critical Control}

\begin{findingbox}
\textbf{No rotation augmentation is applied during training.} All three models see only canonical-orientation shapes. Any accuracy difference under test-time rotation is attributable entirely to the \emph{architecture}, not the training protocol. This directly operationalises van Gemert RP9.
\end{findingbox}

% ─────────────────────────────────────────────────────────────
\section{Results}

All three models were trained for 50 epochs under identical conditions:

\begin{center}
\begin{tabular}{ll}
\toprule
\textbf{Hyperparameter} & \textbf{Value} \\
\midrule
Optimiser & Adam, $lr = 10^{-3}$ \\
Loss & Label-smoothing cross-entropy ($\varepsilon = 0.2$) \\
LR schedule & ReduceLROnPlateau (factor 0.5, patience 5) \\
Epochs & 50 \\
Rotation augmentation & \textbf{None} \\
\bottomrule
\end{tabular}
\end{center}

\begin{figure}[h]
  \centering
  \includegraphics[width=\linewidth]{rotation_accuracy.png}
  \caption{%
    \textit{Left:} Accuracy vs.\ Y-axis rotation angle --- the single controlled variable (van Gemert RP9).
    \textit{Right:} Full SO(3) random rotation, mean $\pm$ std over 3 seeds.
    Dashed line $= 20\%$ random chance.
  }
  \label{fig:results}
\end{figure}

\begin{table}[h]
\centering
\caption{Classification accuracy (\%) per rotation condition. All models trained on aligned data only.}
\label{tab:results}
\begin{tabular}{lcccccccc}
\toprule
\textbf{Model} & $\mathbf{0^\circ}$ & $\mathbf{30^\circ}$ & $\mathbf{60^\circ}$ & $\mathbf{90^\circ}$ & $\mathbf{120^\circ}$ & $\mathbf{150^\circ}$ & $\mathbf{180^\circ}$ & \textbf{SO(3)} \\
\midrule
SimpleView & 100 & 96.6 & 60 & 60 & 60  & 93.2 & 100 & 60.3 \\
PointNet   & 100 & 100  & 80 & 60 & 63.4 & 98  & 80  & \textbf{66.7} \\
DGCNN      & 100 & 80   & 42 & 59.8 & 60 & 76.8 & 80  & 53.4 \\
\midrule
Random     & 20  & 20   & 20 & 20 & 20   & 20   & 20  & 20.0 \\
\bottomrule
\end{tabular}
\end{table}

\subsection*{Van Gemert Storyline (c1--c4)}

\storypoint{c1 — Controlled setting is valid}{%
All three models reach 100\% at $0^\circ$. The shapes are distinguishable and all models are fairly set up. The controlled setting is a reasonable approximation (van Gemert c1). $\checkmark$}

\storypoint{c2 — Baselines are reasonable}{%
PointNet and DGCNN both reach 100\% on aligned data. No model is artificially disadvantaged. $\checkmark$}

\storypoint{c3 — The problem exists}{%
All three methods drop to $\approx 60\%$ at $90^\circ$--$120^\circ$. Rotation clearly hurts classification when models are trained without augmentation. $\checkmark$}

\storypoint{c4 — Does SimpleView address the problem? \textbf{No.}}{%
SimpleView (60.3\% SO(3)) does \textbf{not} outperform PointNet (66.7\% SO(3)). The pre-experiment hypothesis is \textbf{not supported}. In fact, PointNet's explicit T-Net canonicalisation marginally outperforms all others.}

\subsection*{Why Does Accuracy Floor at Exactly 60\%?}

Three of the five shapes are symmetric under Y-axis rotation: \textbf{sphere} (fully symmetric), \textbf{torus} (Y-symmetric), and \textbf{cube} (all faces are identical squares). These are always classified correctly. The \textbf{cylinder} and \textbf{cone} have a preferred axis --- their appearance changes substantially at intermediate angles --- and all models fail on them. Floor $= 3/5 = 60\%$.

\subsection*{Why Doesn't SimpleView Help?}

SimpleView's six cameras point in fixed world-space directions ($\pm X, \pm Y, \pm Z$). When a cylinder rotates $90^\circ$ around Y, the $+Z$ view now shows the cylinder cap instead of its lateral surface --- a view pattern \emph{never seen during aligned training}. ResNet-18 has no mechanism to generalise from this. Crucially, there is \textbf{no pooling over view directions}: each view contributes a separate feature vector that is concatenated. An unseen view produces an unseen feature.

\subsection*{The V-Shaped Accuracy Curve}

Accuracy is high at both $0^\circ$ and $180^\circ$ because a $180^\circ$ Y-rotation maps shapes back to an appearance very similar to $0^\circ$: the cone apex still points up, the cylinder is still vertical. All five classes are classified correctly at $180^\circ$.

% ─────────────────────────────────────────────────────────────
\section{Conclusion}

A synthetic five-class dataset isolated rotation as the single controlled variable to test whether SimpleView's multi-view projection is intrinsically more rotation-robust than PointNet and DGCNN.

\begin{findingbox}
\textbf{All three architectures degrade equally under rotation when trained without augmentation.} PointNet (66.7\% SO(3)) marginally outperforms SimpleView (60.3\%) thanks to its T-Net canonicalisation --- directly contradicting the pre-experiment hypothesis.

\medskip
SimpleView's published advantage on ScanObjectNN is real, but this controlled experiment shows it does \textbf{not} come from rotation robustness. The likely sources are noise, partial occlusion, and background clutter --- factors deliberately excluded here. Testing those would require a separate controlled experiment with a different single variable (van Gemert RP9).
\end{findingbox}

\subsection*{Limitations}

\begin{itemize}[noitemsep, leftmargin=2em]
  \item Only five simple geometric primitives; real-world objects may yield different conclusions.
  \item 50 epochs without rotation augmentation; augmented training might change the relative ranking.
  \item Results based on a single training seed per model; repeated seeds would reduce variance.
\end{itemize}

% ─────────────────────────────────────────────────────────────
\section{Links to Artifacts}

\begin{itemize}[leftmargin=2em]
  \item \textbf{Dataset:} \href{https://github.com/YOUR\_USERNAME/simpleview-rotation-experiment}{[Add Google Drive / Zenodo link]}\\
    11 HDF5 files ($\approx$85 MB): training set + 7 Y-axis rotation conditions + 3 SO(3) seeds.
  \item \textbf{Generation \& training code:} \href{https://github.com/YOUR\_USERNAME/simpleview-rotation-experiment}{github.com/YOUR\_USERNAME/simpleview-rotation-experiment}
\end{itemize}

\subsection*{How to Reproduce}

\begin{codebox}
\# 1. Generate the dataset\\
python code/generate\_dataset.py\\[4pt]
\# 2. Train all three models (requires CUDA GPU)\\
python code/train\_eval.py -{}-models simpleview pointnet dgcnn\\[4pt]
\# 3. Evaluate and combine results\\
python code/eval\_dgcnn\_combine.py\\[4pt]
\# 4. Generate all figures (scripted, RP11)\\
python code/plot\_results.py
\end{codebox}

% ─────────────────────────────────────────────────────────────
\section*{References}

\begin{itemize}[leftmargin=2em, label={}, itemsep=3pt]
  \item Goyal, A. et al.\ (2021). \textit{SimpleView: Revisiting the use of depth maps for 3D object detection}. ICML. \href{https://arxiv.org/abs/2106.05304}{arXiv:2106.05304}
  \item Qi, C.\ R. et al.\ (2017). \textit{PointNet: Deep Learning on Point Sets for 3D Classification and Segmentation}. CVPR. \href{https://arxiv.org/abs/1612.00593}{arXiv:1612.00593}
  \item Wang, Y. et al.\ (2019). \textit{Dynamic Graph CNN for Learning on Point Clouds}. ACM TOG. \href{https://arxiv.org/abs/1801.07829}{arXiv:1801.07829}
  \item van Gemert, J. (2024). \textit{Research Guidelines in Deep Learning}. \href{https://jvgemert.github.io/ResearchGuidelinesInDL.pdf}{jvgemert.github.io}
\end{itemize}

\end{document}
